%!TEX TS-program = xelatex
%!TEX encoding = UTF-8 Unicode
% Awesome CV LaTeX Template for CV/Resume
%
% This template has been downloaded from:
% https://github.com/posquit0/Awesome-CV
%
% Author:
% Claud D. Park <posquit0.bj@gmail.com>
% http://www.posquit0.com
%
% Template license:
% CC BY-SA 4.0 (https://creativecommons.org/licenses/by-sa/4.0/)
%


%-------------------------------------------------------------------------------
% CONFIGURATIONS
%-------------------------------------------------------------------------------
% A4 paper size by default, use 'letterpaper' for US letter
\documentclass[11pt, a4paper]{awesome-cv}

\usepackage{datenumber,fp}

% Configure page margins with geometry
% \geometry{left=1.4cm, top=.8cm, right=1.4cm, bottom=1.8cm, footskip=.5cm}
\geometry{left=1.4cm, top=.75cm, right=1.4cm, bottom=1.1cm, footskip=.5cm}

% Specify the location of the included fonts
\fontdir[fonts/]

% Color for highlights
% Awesome Colors: awesome-emerald, awesome-skyblue, awesome-red, awesome-pink, awesome-orange
%                 awesome-nephritis, awesome-concrete, awesome-darknight
% \colorlet{awesome}{awesome-red}
\definecolor{googleblue}{HTML}{4285F4}
\colorlet{awesome}{googleblue}
% Uncomment if you would like to specify your own color
% \definecolor{awesome}{HTML}{CA63A8}

% Colors for text
% Uncomment if you would like to specify your own color
% \definecolor{darktext}{HTML}{414141}
% \definecolor{text}{HTML}{333333}
% \definecolor{graytext}{HTML}{5D5D5D}
% \definecolor{lighttext}{HTML}{999999}

% Set false if you don't want to highlight section with awesome color
\setbool{acvSectionColorHighlight}{true}

% If you would like to change the social information separator from a pipe (|) to something else
\renewcommand{\acvHeaderSocialSep}{\quad\textbar\quad}


%-------------------------------------------------------------------------------
%	PERSONAL INFORMATION
%	Comment any of the lines below if they are not required
%-------------------------------------------------------------------------------
% Available options: circle|rectangle,edge/noedge,left/right
% \photo[rectangle,edge,right]{./examples/profile}
\name{Kai-Hung}{Huang}
% \position{Sr. Software Engineer at Xperi (DTS){\enskip\cdotp\enskip}Android Expert}
\position{Senior Software Engineer \enskip\cdotp\enskip Android Platform & Embedded Systems}
% \address{3F., No. 3, Ln. 180, Sec. 1, Zhongshan N. Rd., Tamsui Dist., New Taipei City 251017, TAIWAN}
% \address{5F., No. 65, Ln. 182, Beixin Rd., Tamsui Dist., New Taipei City 251608, Taiwan}
% \address{7F., No. 262-5, Sec. 1, Xinshi 1st Rd., Tamsui Dist., New Taipei City 251652, Taiwan}
\address{Taipei, Taiwan}
\mobile{(+886) 0925-133902}
\email{randy.hunang@gmail.com}
% \github{Randy-Huang}
\linkedin{kai-hung-huang-344117106}
% \gitlab{gitlab-id}
% \stackoverflow{SO-id}{SO-name}
% \twitter{@twit}
% \skype{skype-id}
% \reddit{reddit-id}
% \extrainfo{extra informations}

% \quote{``Be the change that you want to see in the world."}

\makeatletter
\renewcommand*{\skilltypestyle}[1]{{\fontsize{9.5pt}{1em}\bodyfont\mdseries\color{darktext} #1}}
\makeatother

%-------------------------------------------------------------------------------
\begin{document}

% Print the header with above personal informations
% Give optional argument to change alignment(C: center, L: left, R: right)
\makecvheader[C]

% Print the footer with 3 arguments(<left>, <center>, <right>)
% Leave any of these blank if they are not needed
% \makecvfooter
%  {\today}
%  {Kai-Hung Huang~~~·~~~Résumé}
%  {\thepage}
\makecvfooter{}{Kai-Hung Huang \enskip\cdotp\enskip R\'esum\'e}{\thepage}



%-------------------------------------------------------------------------------
%	CV/RESUME CONTENT
%	Each section is imported separately, open each file in turn to modify content
%-------------------------------------------------------------------------------
%-------------------------------------------------------------------------------
%	SECTION TITLE
%-------------------------------------------------------------------------------
\cvsection{Summary}

%-------------------------------------------------------------------------------
%	SUBSECTION TITLE
%-------------------------------------------------------------------------------
\cvsubsection{Objective}

%-------------------------------------------------------------------------------
%	CONTENT
%-------------------------------------------------------------------------------
\begin{cvparagraph}

%---------------------------------------------------------
    Seeking an Android software development,especially in Framewrk/API/Application layers, where I specialized in Java/C/C++ and 7+ work on AOSP.

\end{cvparagraph}

%-------------------------------------------------------------------------------
%	SUBSECTION TITLE
%-------------------------------------------------------------------------------
\cvsubsection{About}

%-------------------------------------------------------------------------------
%	CONTENT
%-------------------------------------------------------------------------------
\begin{cventries}

%---------------------------------------------------------
  \cvsingleentry
    {
      \begin{cvitems} % Description(s)
        \item {9+ years experienced embedded Android/Linux system software (Android Framework/API/Application layer).}
        \item {4+ years specialized in Android Framework (JNI/SystemService/Bluetooth).}
        \item {3+ years developed apps with MVVM architecture, adept at utilizing DI and using LiveData and ViewModel for data management.}
        \item {2+ years supervised a team to efficiently customize SurfaceFlinger/API/Settings for integrating features to 3rd developer.}
        \item {\textasciitilde1 year worked on sensor domain of Android BSP development.}
        \item {A passionate person who focuses on programming and has good communication skills, and excellent team player.}
      \end{cvitems}
    }

\end{cventries}

%-------------------------------------------------------------------------------
% SECTION TITLE
%-------------------------------------------------------------------------------
% \cvsection{Technical Skills}

%-------------------------------------------------------------------------------
% CONTENT
%-------------------------------------------------------------------------------
%\begin{cvskills}
%  \cvskill{Languages}{C/C++, Java, Kotlin}
%  \cvskill{Android \& Linux}{AOSP, Android Framework, JNI, System Services, Android NDK, Linux system programming}
%  \cvskill{Concurrency \& IPC}{concurrent programming and thread synchronization, Unix domain sockets, and shared-memory ring buffers for audio/camera data transfer}
%  \cvskill{Connectivity}{Wi-Fi / SoftAP, Bluetooth Classic / BLE, RTP/RTCP Data Streaming}
%  \cvskill{Platform Integration}{Android HAL/BSP, SEPolicy, sensors, Linux driver integration}
%  \cvskill{Performance \& Debugging}{Systrace, Android Studio Profiler, CPU and memory analysis, system stability debugging}
%\end{cvskills}

\cvsection{Technical Skills}

\begin{cvskills}
  \cvskill{Languages}{C/C++, Java, Kotlin, Python, Shell Script}
  \cvskill{Android OS \& Kernel}{AOSP, Android Framework (JNI, System Services, SurfaceFlinger), Android NDK, HAL, Linux Systems Programming}
  \cvskill{Protocols \& IPC}{Bluetooth (Classic/BLE), Wi-Fi/SoftAP, RTP/RTCP, Unix Domain Sockets, Ashmem}
  \cvskill{Tools \& Debugging}{Systrace, Android Studio Profiler, Perfetto, Bugreport Analysis, Git, Audacity}
\end{cvskills}


%-------------------------------------------------------------------------------
%	SECTION TITLE
%-------------------------------------------------------------------------------
\cvsection{Work Experience}

%-------------------------------------------------------------------------------
%	CONTENT
%-------------------------------------------------------------------------------
\begin{cventries}

%---------------------------------------------------------
  \cventry
    {Sr. Software Engineer} % Job title
    {XPERI Corporation (DTS)} % Organization
    {Taipei, Taiwan} % Location
    {Aug. 2021 - Present} % Date(s)
    {
      \begin{cvitems} % Description(s) of tasks/responsibilities
        \item {Developed list-style layout using the extended RecyclerView with GridLayoutManager and customized Adapter for optimized UI representation.}
        \item {Maintained streaming techiques with RTP/RTCP network protocols transmitting audio data for DTS Play-Fi wireless multi channel.}
        \item {Customized a linux process which can receive audio PCM data from system using unix domain socket based on Android TV.}
      \end{cvitems}
    }

%---------------------------------------------------------
  \cventry
    {Lead Programmer} % Job title
    {XRSPACE Inc.} % Organization
    {Taipei, Taiwan} % Location
    {Mar. 2019 - Aug. 2021} % Date(s)
    {
      \begin{cvitems} % Description(s) of tasks/responsibilities
        \item {Lead team members (~4 people) to reach the project tasks.}
        \item {Coordinated and host task topics for cross-team collaboration to decrease unnecessary effort.}
        \item {Developed a library that connects between HMD and mobile device using BT module, which customized encryption/cryption and packet.}
        \item {Customized SurfaceView/SurfaceFlinger to cast specified layers of VR content from a foreground VR applicatoin to TV by Miracast.}
        \item {Implemented a library dedicating specified CPU to a particular process or function by using the technique of CPU affinity.}
        \item {Experienced with analyzing the system performance and stability using cmd tools/systrace/Android profiler.}
        \item {Experienced with fine-tuning an application performance for UI rendering with systrace.}
      \end{cvitems}
    }

%---------------------------------------------------------
  \cventry
    {Sr. Software Engineer} % Job title
    {} % Organization
    {} % Location
    {Sep. 2018 - Mar. 2019} % Date(s)
    {
      \begin{cvitems} % Description(s) of tasks/responsibilities
        \item {Developed camera (mono, RGB, depth) getting from QVR (Qualcomm VR solution) and Android Camera NDK for VR 6-Dof processing.}
        \item {Implemented a Ring Buffer data structure to store camera frame and transmit it to clients (Java/Unity layers) through ashmem or fd.}
        \item {Designed a preview window tool by ANativeWindow to display camera frame using YUV-420/10-bit Raw RGB format.}
        \item {Designed communication architecture between VR headset and Android phone using Bluetooth LE.}
      \end{cvitems}
    }

%---------------------------------------------------------
  \cventry
    {Sr. Software Engineer} % Job title
    {ASUSTeK Computer Inc.} % Organization
    {Taipei, Taiwan} % Location
    {Dec. 2015 - Sep. 2018} % Date(s)
    {
      \begin{cvitems} % Description(s) of tasks/responsibilities
        \item {[Zenbo Junior]Developed the behavior of "Line Following" feature with PID control algorithm.}
        \item {[Zenbo]Developed RobotAPI and RobotSDK for 3rd Developer based on Android Framework.}
        \item {[Zenbo]Developed the behavior of moving to Docking charger and integrate MCU command for charger control.}
        \item {[Zenbo]Integrated Battery/Facial/Motion features of main functions using Android Service to RobotFramework.}
        \item {[Zenbo]Android BSP development - integrate motion sensors (accelerometer, gyroscope, and compass) using Intel ISH (sensor hub).}
        \item {[Zenbo]Android BSP development - porting capacitive touch sensor (GPIO, BIOS, Linux driver, Android HAL/Framework).}
        \item {Android App unit test tool for sensor testing/debugging.}
      \end{cvitems}
    }

    %---------------------------------------------------------
\end{cventries}

%\input{resume/honors.tex}
% \input{resume/presentation.tex}
%\input{resume/writing.tex}
% \input{resume/committees.tex}
%-------------------------------------------------------------------------------
%	SECTION TITLE
%-------------------------------------------------------------------------------
\cvsection{Education}


%-------------------------------------------------------------------------------
%	CONTENT
%-------------------------------------------------------------------------------
\begin{cventries}

%---------------------------------------------------------
  \cventry
    {M.Eng., Robotics Engineering} % Degree
    {Tamkang University} % Institution
    {New Taipei City, Taiwan} % Location
    {Mar. 2013 - Aug. 2014} % Date(s)
    {
      \begin{cvitems} % Description(s) bullet points
        % \item {Grade: GPA: 4 and Course Average: 91.36/100.}
        % \item {Graduated in one academic year (Excellent MS student qualification)}
        % \item {Scholarships (over \$100,000+): Awarded outstanding academic achievement for 2 semesters}
	\item {GPA: 4.0/4.0; course average: 91.36/100. Completed the program in one academic year.}
      \end{cvitems}
    }

%---------------------------------------------------------
  \cventry
    {B.Eng., Electrical Engineering} % Degree
    {Tamkang University} % Institution
    {New Taipei City, Taiwan} % Location
    {Mar. 2009 - Aug. 2013} % Date(s)
    {
      \begin{cvitems} % Description(s) bullet points
        % \item {Grade: GPA: 4 and Course Average: 92.60/100.}
        % \item {Academic Honors: Graduated in the honor of the top-rated prize, and Awarded 1st excellent academic for 7 semesters.}
        % \item {Scholarships (over \$100,000+): Awarded excellent academic for 8 consecutive semesters.}
	% \item {GPA: 4.0/4.0; course average: 92.60/100.}
	\item {GPA: 4.0/4.0; graduated first in class; recipient of Academic Excellence Scholarships for 8 consecutive semesters.}
      \end{cvitems}
    }

%---------------------------------------------------------
\end{cventries}

%\input{resume/extracurricular.tex}


%-------------------------------------------------------------------------------
\end{document}
